\hypertarget{appendix-ed338677de0a}{}\label{referencestart:appendix-ed338677de0a}
\section{Checksums, CRCs, and hashes}\label{appendix-ed338677de0a-checksums-crcs-and-hashes}

\textbf{Curriculum depth:} Reference implemented and known-answer tested · \textbf{First seen:} C006 · \textbf{Latest development:} Phase 28 repair

A checksum or digest maps an arbitrary byte sequence to a smaller fixed-size value. Different families make different promises.

An additive checksum is mainly an accidental-error detector and is weak against structured changes such as reordering or compensating byte modifications.

A CRC interprets bits as a polynomial over \(GF(2)\) and computes a remainder under a declared generator polynomial and parameter set. CRCs are excellent for many transmission/storage error patterns, but they are not cryptographic authenticators.

A cryptographic hash such as SHA-256 is designed so that finding collisions or preimages is computationally hard under the security assumptions of the primitive. A hash still does \textbf{not} authenticate a file when an attacker can replace both the file and the untrusted hash value.

\subsection{SHA-256 compression structure}\label{appendix-ed338677de0a-sha-256-compression-structure}

SHA-256 pads the message with a one bit, enough zero bits to leave 64 bits at the end of the final 512-bit block, and then the original message length modulo \(2^{64}\). Each 512-bit block expands sixteen 32-bit words into a 64-word schedule and updates an eight-word chaining state through 64 rounds.

The important learning boundary is representation: message words are interpreted in the algorithm's specified byte order. Reinterpreting native host integers without explicit conversion can produce a platform-specific hash implementation.

\subsection{Known-answer evidence}\label{appendix-ed338677de0a-known-answer-evidence}

tools/labs/crypto\_primitives\_lab.cpp implements SHA-256 directly and validates the standard digest of ``abc'' and the empty message. The same SHA-256 implementation is then used by HMAC, HKDF and a domain-separated Merkle tree.

This raises the evidence above ``definition present,'' but the implementation is educational. It is not reviewed for constant-time behavior, hardware acceleration, fault attacks, side channels or production cryptographic use.

\subsection{What to use for what}\label{appendix-ed338677de0a-what-to-use-for-what}

Use a checksum/CRC when the threat is accidental corruption and the protocol specifies the polynomial/parameters.

Use a cryptographic hash when content identity or collision/preimage resistance is required.

Use a MAC such as HMAC when a party holding a secret key must authenticate content.

Use a digital signature when verification must be possible without sharing the signing secret.

\subsection{Prerequisites}\label{appendix-ed338677de0a-prerequisites}

\hyperref[appendix-d49161831368]{Bitwise XOR}, \hyperref[appendix-91d09ebed979]{Serialization}, \hyperref[appendix-bc07c3d25b18]{Endianness}.

\subsection{Used by}\label{appendix-ed338677de0a-used-by}

\hyperref[appendix-cbe594d1ab0a]{Integrity and authenticity}, \hyperref[appendix-50af792c9da7]{HMAC and HKDF}, \hyperref[appendix-098192f8d46d]{Merkle trees and authenticated structure}, \href{../07-Vision-XR-and-Playable-Worlds/main.pdf\#nameddest=appendix-91eaa3cbdbc3}{Production manifests and transactional publication (Vision, XR, and Playable Worlds, Reference notes: Production manifests and transactional publication)}.

\subsection{Source and navigation}\label{appendix-ed338677de0a-source-and-navigation}

\hyperref[01-foundations-i-ch002-bytes-become-a-durable-agreement]{Continuous teaching story} · \href{../07-Vision-XR-and-Playable-Worlds/main.pdf\#nameddest=appendix-4ad47bdd269a}{Dependency Map (Vision, XR, and Playable Worlds, Reference notes: Dependency Map)}
